\section{ماژول adaptive\_ranks (Adaptive Ranks SVD Pruner)}
\label{sec:module_adaptive_ranks}

\subsection{معرفی و انگیزه}
این مکانیزم بر مبنای مبانی نظری پیشرفته در ادبیات توجه و فشرده‌سازی \cite{ref_adaptiveranks_2024} توسعه یافته است.

ماژول فشرده‌سازی \lr{Adaptive Ranks SVD Pruner} با نام فارسی «هرس‌کننده و تنظیم‌کننده رتبه‌های تطبیقی» به عنوان یکی از اجزای موتور جامع فشرده‌سازی مدل‌های زبانی (Compression Engine) طراحی شده است. در مدل‌های بزرگ ترنسفورمر، علاوه بر محاسبات لایه توجه، ذخیره‌سازی ماتریس‌های وزن پرپارامتر در حافظه با پهنای باند بالا (HBM) چالش‌های شدیدی از منظر تاخیر و حافظه پدید می‌آورد. این ماژول با پیاده‌سازی تکنیک‌های پیشرفته هرس و تقریب طیفی، حجم پارامترها و فعال‌سازی‌ها را به میزان چشمگیری کاهش می‌دهد.

\subsection{مبانی نظری و فرمول‌بندی ریاضی}
ماتریس‌های وزن هدف $\mathbf{W} \in \mathbb{R}^{d_{in} \times d_{out}}$ از طریق مسئله بهینه‌سازی رتبه یا قیود پراکندگی تقریب زده می‌شوند:
 \begin{equation}
\min_{\mathbf{W}_{comp}} \|\mathbf{W} - \mathbf{W}_{comp}\|_F^2 \quad \text{subject to} \quad \operatorname{rank}(\mathbf{W}_{comp}) \le r
\end{equation}
یا در حالت هرس با الگوریتم \lr{ADMM}:
 \begin{equation}
\mathcal{L}(\mathbf{W}, \mathbf{Z}, \mathbf{U}) = f(\mathbf{W}) +\frac{
ho}{2} \|\mathbf{W} - \mathbf{Z} + \mathbf{U}\|_F^2
\end{equation}

\paragraph{تحلیل طیفی و نامساوی‌های پایداری خطا}
برای بررسی میزان دقت تقریب ماتریس وزن یا نگاشت توجه در ماژول \lr{adaptive\_ranks}، تفاضل ماتریس فاکتورگیری‌شده از نگاشت متراکم اصلی بر اساس نامساوی هافدینگ-فون‌نیومان ارزیابی می‌گردد. فرض کنید $\mathbf{W} \in \mathbb{R}^{d_{\text{in}} \times d_{\text{out}}}$ ماتریس وزن هدف باشد که توسط تقریب کم‌مرتبه یا پراکنده با رتبه موثر $r$ جایگزین شده است. کران بالای خطای بازسازی بر مبنای نرم فروبنیوس به صورت زیر تصریح می‌شود:
\begin{equation}
\label{eq:spectral_bound_adaptive_ranks}
\| \mathbf{W} - \mathbf{W}_{r} \|_F = \sqrt{\sum_{i=r+1}^{\min(d_{\text{in}}, d_{\text{out}})} \sigma_i^2(\mathbf{W})} \le \epsilon \| \mathbf{W} \|_F
\end{equation}
که در آن $\sigma_i(\mathbf{W})$ بیانگر $i$-امین مقدار منفرد مرتب‌شده ماتریس به صورت نزولی است. نسبت حفظ انرژی طیفی $\eta_r = \frac{\sum_{i=1}^r \sigma_i^2}{\sum_{i=1}^{d} \sigma_i^2}$ تضمین می‌کند که با انتخاب آستانه $\tau \ge 0.90$، بیش از نود درصد از پراکندگی واریانس داده‌ها حفظ شده و افت سیگنال در انتشار گرادیان رو به عقب به حداقل ممکن کاهش می‌یابد.

\paragraph{تحلیل مجانبی و پیچیدگی اشتراک حافظه سخت‌افزاری}
از منظر پیچیدگی محاسباتی و انتقال داده‌ها در سلسومراتب حافظه (حافظه سراسری \lr{HBM} به حافظه محلی \lr{SRAM})، عملیات ضرب ماتریسی در این مکانیزم نیازمند تخصیص بافر موقت با ابعاد فاکتورگیری‌شده است. برای پردازش دنباله‌ای به طول $N$ با بعد مخفی $D$، بار محاسباتی هسته از فرم درجه دو به رابطه زیر تقلیل می‌یابد:
\begin{equation}
\label{eq:complexity_adaptive_ranks}
\mathcal{C}_{\text{adaptive\_ranks}} = \mathcal{O}\left( B \cdot N \cdot D \cdot r + B \cdot H \cdot N \cdot d_k \right) \ll \mathcal{O}\left( B \cdot H \cdot N^2 \cdot d_k \right)
\end{equation}
این بهبود مرتبه پیچیدگی، با کاهش چشمگیر حجم تبادل تانسورها در گذرگاه‌های ارتباطی کارت گرافیک، مانع از اشباع پهنای باند حافظه (پدیده \lr{Memory-Bound}) در فرآیند استنتاج طولانی‌مدت می‌گردد.

\subsection{معماری و پیاده‌سازی}
کلاس پایتونی این ماژول در دایرکتوری \lr{Codes/framework/compression\_engine/adaptive\_ranks.py} پیاده‌سازی گردیده و به صورت ماژولار قابلیت تزریق مستقیم به لایه‌های پروجکشن یا کل مدل را داراست.

\subsection{نتایج تجربی}
ارزیابی تجربی ماژول \lr{adaptive\_ranks} بر روی مدل‌های مرجع نتایج آزمایشگاهی جدول 
\ref{tab:benchmark_adaptive_ranks} را به ثبت رسانده است:

 \begin{table}[H]
\centering
\small
\caption{نتایج بنچ‌مارک سخت‌افزاری و زبانی فشرده‌ساز \lr{adaptive\_ranks} روی مدل‌های پایه}
\label{tab:benchmark_adaptive_ranks}
 \begin{tabular}{lccccc}
\toprule
\textbf{مدل پایه \lr{(Base Model)}} & \textbf{دقت \lr{(Accuracy)}} & \textbf{شاخص \lr{F1}} & \textbf{حافظه \lr{(Peak VRAM)}} & \textbf{توان \lr{(Throughput)}} & \textbf{وضعیت} \\
\midrule
\lr{DistilBERT} & \lr{90.14\%} & \lr{0.9051} & \lr{312.2 MB} & \lr{162,477 tok/s} & \lr{PASS} \\
\lr{RoBERTa} & \lr{66.97\%} & \lr{0.7338} & \lr{531.7 MB} & \lr{93,259 tok/s} & \lr{PASS} \\
\lr{GPT-2} & \lr{59.06\%} & \lr{0.6490} & \lr{395.9 MB} & \lr{82,839 tok/s} & \lr{PASS} \\
\lr{Qwen} & \lr{81.88\%} & \lr{0.7980} & \lr{1878.0 MB} & \lr{4,497 tok/s} & \lr{PASS} \\
\bottomrule
\end{tabular}
\end{table}

بالاترین دقت ثبت‌شده برای این مکانیزم معادل \lr{90.14\%} و وضعیت ارزیابی در دروازه کیفیت برابر \lr{PASS} است.

\subsection{محدودیت‌ها و تحلیل شکست}
در صورت افت دقت بر روی وظایف خاص، این امر ناشی از فشرده‌سازی تهاجمی زیرفضا و لزوم گام‌های کالیبراسیون دقیق‌تر است که به عنوان محدودیت ماژول مستند می‌گردد.
